The definitions are in \path{TaylorDiscriminants/ConnectionKillsSeries.lean}, \path{nabla}, lines 31--34, \path{Hset}, lines 41--43, and \path{ConnectionIntegrality}, lines 77--83. Under an integral formal biextension, a canonical gauge, and $w_0+w_0^*=1$, the supplied theorem
\path{TaylorDiscriminants/WickVanishing/ExpConjConnectionIntegral.lean},
\path{exp_conj_connection_integral}, lines 40--82, proves
\[
 e^{-kh_L}\nabla_{u,\phi-Su}(e^{kh_L}G)\in R[[X,Z]]
 \quad(u,\phi\in R^d,\ G\in R[[X,Z]]).
\]
Its \path{connection_first_identity}, lines 84--90, packages that conclusion. Thus full-series nonintegrality does not by itself obstruct the chosen connection operations.

The supplied ordered-column theorem has a denominator budget rather than a height-dependent positive slope. In \path{TaylorDiscriminants/GermSC.lean}, \path{SCSt}, lines 29--41, and \path{GermSCCols.lean}, \path{SCColsSt}, lines 10--22, take a fractional submodule $J\subset K$ containing one, lattice directions in $H_h$, correction entries in $J$, and the displayed vanishing of the three symplectic pairings except for equal colors and opposite sides. If an integral polynomial has columnwise total \path{pairMin} at most $e$, its ordered connection value on $e^{kh_L}f$, for integral $f$, lies in $J^e$. These are supplied by \path{WickVanishing/Closed.lean}, \path{symmetric_canonical_gauge}, lines 25--26, and \path{scaled_connection_columns}, lines 47--63, with \path{lemma_H_columns}, lines 22--23, providing the stated selected-derivative comparison. For $J=p^{-b}R$, the conclusion is a lower valuation $-be$. A new positive higher-height estimate for the actual lattice, corrections and selected operators is additional quantitative work.

Canonical-gauge independence is also supplied, by
\path{Development/AbelianVarietyTheory/Proofs/FormalGroupCharts/CanonGaugeIndependenceAll.lean},
\path{GlobalTheory.canonGauge_independence_all}, lines 23--26. The difference of two canonical gauges is balanced and biadditive; its diagonal in logarithm coordinates is quadratic. The selected value is unchanged when the associated second-order operator annihilates $P$. This removes the specified quadratic ambiguity. It does not remove an arbitrary higher-degree primitive kernel: the series \path{be} in \path{ProductZerothFirstParts.lean}, lines 186--201, is an arbitrary multivariable series. DP38--DP39 below give exact tests of that distinction. The complete source-read scope, including which of the recovered modules were actually read, is recorded in \path{research/theta_normalization_harmonic_trace.md}.

